
\subsubsection{Interpreting the Existence Result}

The 99.51\% fingerprint classification accuracy demonstrates that fine-tuning creates highly discriminative representation signatures. The effect size (Cohen's d = 698.08) indicates near-complete separation between benchmark-specific representation spaces. The asymmetric confusion matrix---CIFAR-100 features never misclassified while 49 Flowers102 samples were---suggests different benchmarks create fingerprints of varying strength or distinctiveness.

This finding is consistent with the hypothesis that fine-tuning encodes benchmark-specific statistical regularities. Flowers102 contains fine-grained petal textures and natural backgrounds; CIFAR-100 contains diverse object categories at low resolution. These systematic differences produce linearly separable representations.

\subsubsection{Why BFS Failed to Predict Gap}

Several factors may explain the null correlation:

\begin{enumerate}
\item \textbf{Ceiling effect:} All models achieved BFS > 0.999. With effectively zero variance in BFS, correlation is mathematically impossible. This is a metric saturation problem, not necessarily evidence against an underlying relationship.
\end{enumerate}

\begin{enumerate}
\item \textbf{Binary fingerprints:} The fingerprint phenomenon may be presence/absence rather than graded. Once representations encode benchmark-specific patterns, additional ''fingerprint intensity'' may not exist.
\end{enumerate}

\begin{enumerate}
\item \textbf{Domain shift dominance:} CIFAR-100 (coarse-grained objects) and Flowers102 (fine-grained plants) represent fundamentally different visual domains. The observed gap (33-44 percentage points) may primarily reflect domain distance rather than benchmark-specific overfitting.
\end{enumerate}

\begin{enumerate}
\item \textbf{Sample size limitation:} With n = 6 models, statistical power for detecting moderate correlations is low. A true effect could exist but remain undetected.
\end{enumerate}

\subsubsection{Limitations}

\begin{enumerate}
\item \textbf{Reduced benchmark count:} The proof-of-concept used 2 benchmarks instead of the planned 5, increasing chance level from 20\% to 50\% and reducing the stringency of the existence test.
\end{enumerate}

\begin{enumerate}
\item \textbf{BFS metric saturation:} The proposed BFS metric saturates at high values, preventing correlation analysis. Alternative metrics (entropy, margin, calibrated probabilities) may provide more variance.
\end{enumerate}

\begin{enumerate}
\item \textbf{H-M2 not executed with real data:} The single vs. multi-benchmark comparison remains untested with real datasets.
\end{enumerate}

\begin{enumerate}
\item \textbf{Single architecture:} Only ResNet-50 was tested. Vision Transformers and other architectures may exhibit different fingerprint characteristics.
\end{enumerate}

\begin{enumerate}
\item \textbf{Domain heterogeneity confound:} Using CIFAR-100 and Flowers102---domains with substantial differences---conflates domain shift with benchmark-specific effects. Testing within a single domain (e.g., multiple fine-grained datasets) would isolate benchmark effects.
\end{enumerate}

